This chapter documents the application of Parameter-Efficient Fine-Tuning (PEFT) algorithms to enable specialized visual reasoning on localized hardware. Specifically, this research details the Knowledge Distillation sequence utilizing Low-Rank Adaptation (LoRA) and optimized CUDA kernels (Unsloth) to transfer high-level spatial reasoning from a teacher model (GPT-4o) into an 11-Billion parameter student model (Llama-Vision 11B). This process results in a model capable of executing complex routing and layout verification tasks within the 24 GB VRAM envelope of a single consumer Graphics Processing Unit (GPU).

\section{The Edge-Deployment VRAM Bottleneck}
\label{sec:vram_limits}

However, the localized deployment target for this thesis pipeline is an NVIDIA GeForce RTX 3090 possessing 24 GB of VRAM. This hardware restriction necessitates memory optimization strategies. The baseline inference utilizes 4-bit quantization (bitsandbytes), which compresses the model weights while preserving the majority of the original semantic precision.

While 4-bit quantization enables execution, \textit{training} a model conventionally requires storing high-precision gradient and optimizer states. To specialize the Llama-Vision 11B student model on industrial technical layouts without exceeding memory limits, the pipeline employs a combination of LoRA and memory-mapped kernels.

\section{Low-Rank Adaptation (LoRA)}
\label{sec:lora_math}

To train a robust visual referee exclusively on 24 GB of VRAM, this thesis implements Low-Rank Adaptation (LoRA). Rather than explicitly altering the $d \times d$ pre-trained weight matrix $W_0 \in \mathbb{R}^{d \times d}$ during gradient descent, LoRA mathematically freezes $W_0$ entirely. 

The algorithm then introduces two trainable decomposition matrices $A$ and $B$, possessing a constrained rank dimensionality $r$ where $r \ll d$. The forward pass incorporates the low-rank update $\Delta W$ alongside the frozen base weights:

\begin{equation}
h = W_0 x + \Delta W x = W_0 x + BA x
\end{equation}

where $B \in \mathbb{R}^{d \times r}$ and $A \in \mathbb{R}^{r \times d}$. By restricting the rank to $r = 16$, the number of trainable parameters is reduced by over 98\%, significantly lowering the VRAM required for gradient storage while maintaining the model's ability to learn domain-specific visual features.

To train the student model for structured extraction and layout verification, a high-quality spatial dataset was synthesized using a teacher-student paradigm. In this framework, the outputs of a more capable \textit{teacher} model (GPT-4o) serve as the ground truth for training the smaller student.

A dataset of approximately 1,000 industrial engineering documents was processed through the DocLayout-YOLO detector. From each document, various regions (tables, figures, and technical paragraphs) were cropped and passed to the teacher model. The teacher was prompted to generate structured XML and JSON representations of the content, specifically focusing on complex multi-column parameters and schematic descriptions.

These input (image crop) and output (structured text) pairs were then formatted into the multimodal conversation structure required for fine-tuning.
\begin{verbatim}
[
  {
    "role": "user",
    "content": [
      {"type": "image"},
      {"type": "text", "text": "Extract XML Grid <ROW><P><C><V>"}
    ]
  },
  {
    "role": "assistant",
    "content": [
      {"type": "text", "text": "<ROW><P>Voltage</P><C>Max</C><V>24 V</V></ROW>"}
    ]
  }
]
\end{verbatim}

\section{Fine-Tuning Execution and Hyperparameters}
\label{sec:unsloth_training}

To maximize memory efficiency, the training pipeline was implemented using the \textit{Unsloth} framework, which provides optimized kernels for LoRA fine-tuning. The student model was loaded in 4-bit precision, and the LoRA adapters were applied to the $q, k, v, o$ and $gate, up, down$ projections within the transformer blocks.

The training followed a precise hyperparameter schedule: a learning rate of $2 \times 10^{-4}$ with a cosine decay, a rank $r=16$, and an alpha $\alpha=32$. Training was executed for 1,000 steps with a global batch size of 8 (achieved through gradient accumulation), taking approximately 4 hours on the target RTX 3090 hardware.

Figure \ref{fig:training_loss} shows the training loss trajectory logged via Weights \& Biases. The loss curve demonstrates smooth convergence, indicating that the model successfully learned the structured extraction task without overfitting to the limited variety of the training set. 

Upon completion, the LoRA adapters were merged with the base model and pushed to a private repository (`RMunshi/vlm-student-thesis`) for integration into the production pipeline. The resulting specialist model provides the visual reasoning necessary for the routing and verification layers while maintaining the low latency required for industrial document processing.
